% AstreaView - Documentacion del Proyecto Grupal 2 (CE4303), formato IEEE-Trans.
% Overleaf: New Project > Blank Project, pegar todo en main.tex y compilar con pdfLaTeX.
% Todo lo marcado con \pendiente{...} (en rojo) hay que completarlo antes de entregar.
\documentclass[journal]{IEEEtran}

\usepackage[utf8]{inputenc}
\usepackage[T1]{fontenc}
\usepackage[spanish,es-noshorthands,es-tabla]{babel}
\usepackage{amsmath}
\usepackage{graphicx}
\usepackage{booktabs}
\usepackage{tabularx}
\usepackage{xcolor}
\usepackage{listings}
\usepackage{cite}
\usepackage{url}
\usepackage{tikz}
\usetikzlibrary{arrows.meta,positioning,fit,backgrounds}
\usepackage[hidelinks]{hyperref}

\newcommand{\pendiente}[1]{\textcolor{red}{[PENDIENTE: #1]}}

\lstset{
  basicstyle=\ttfamily\scriptsize,
  breaklines=true,
  frame=single,
  columns=fullflexible,
  keepspaces=true
}

\tikzset{
  caja/.style={draw, rounded corners=2pt, align=center, font=\footnotesize,
               minimum height=7mm, inner sep=3pt, fill=white},
  grupo/.style={draw, dashed, rounded corners=4pt, inner sep=6pt, fill=black!4},
  flecha/.style={-{Stealth[length=2mm]}, thick},
  doble/.style={{Stealth[length=2mm]}-{Stealth[length=2mm]}, thick},
  etiqueta/.style={font=\scriptsize, fill=white, inner sep=1pt, align=center}
}

\begin{document}

\title{AstreaView: escaneo 3D distribuido con visualización holográfica}

\author{Emerson~Monge~H.,
        Emanuel~\pendiente{apellido},
        Gabriel~\pendiente{nombre completo y confirmar}%
\thanks{Proyecto Grupal 2, curso Principios de Sistemas Operativos (CE4303),
Ingeniería en Computadores, Tecnológico de Costa Rica \pendiente{confirmar datos,
profesor, semestre y carnés}.}}

\markboth{Principios de Sistemas Operativos -- Proyecto Grupal 2}{AstreaView}
\maketitle

\begin{abstract}
AstreaView es un sistema que escanea un objeto con la cámara TrueDepth de un
iPhone, reconstruye su geometría repartiendo el cálculo entre varios nodos y
muestra el resultado en una computadora y en una pantalla de tipo
\emph{Pepper's Ghost}. El proyecto integra un módulo de kernel propio para la
lectura de botones por GPIO, una reconstrucción volumétrica TSDF distribuida con
una fusión equivalente a la de un solo nodo y una red cifrada con TLS~1.3 y
autenticación mutua. Este documento describe el diseño, el ambiente de
desarrollo, el uso del sistema y los resultados obtenidos hasta la fecha.
\end{abstract}

\begin{IEEEkeywords}
sistemas operativos, controlador de dispositivo, GPIO, procesamiento
distribuido, TSDF, reconstrucción 3D, TLS 1.3, Raspberry Pi.
\end{IEEEkeywords}

% =====================================================================
\section{Introducción}

\IEEEPARstart{U}{n} sistema operativo media entre las aplicaciones y el hardware:
administra interrupciones, memoria, concurrencia y comunicación entre procesos y
equipos. El objetivo del proyecto es implementar las capas necesarias para que
una aplicación interactúe con hardware a través del sistema operativo, cubriendo
cuatro áreas: un controlador propio de entrada y salida, un componente de
hardware, procesamiento distribuido en al menos tres nodos y cifrado de toda la
información que viaja por la red.

AstreaView responde a esas áreas con un escáner 3D de mesa. El objeto se coloca
sobre un plato que se gira a mano y que tiene en el borde un anillo de
marcadores ArUco. Un iPhone transmite imágenes de color y profundidad
(RGB-D), una computadora calcula la pose de cada vista, tres nodos fusionan las
vistas en un volumen común y el modelo resultante se visualiza en la
computadora y en un holograma controlado desde una Raspberry~Pi~3.

\subsection{Controladores y manejo de interrupciones}
En Linux un controlador puede cargarse como módulo de kernel y exponer un
dispositivo de caracteres en \texttt{/dev}, que los procesos usan con las
llamadas \texttt{open}, \texttt{read}, \texttt{poll} e \texttt{ioctl}
\cite{ldd3}. El hardware avisa de sus cambios con interrupciones, cuyo manejador
debe ser breve. Los botones mecánicos producen rebotes, por lo que se necesita un
filtro temporal (\emph{debounce}), implementado aquí con temporizadores de alta
resolución (\texttt{hrtimer}).

\subsection{Reconstrucción volumétrica}
Una función de distancia con signo truncada (TSDF) guarda en cada vóxel la
distancia a la superficie más cercana y un peso de confianza
\cite{curless,kinectfusion}. Al ser un promedio ponderado, varios volúmenes
parciales $d_i$ con pesos $w_i$ se combinan como
\begin{equation}
  D = \frac{\sum_i w_i d_i}{\sum_i w_i}, \qquad W = \sum_i w_i ,
\end{equation}
lo que permite repartir las vistas entre nodos. La superficie se extrae con
\emph{marching cubes} \cite{marching}, se suaviza con el filtro de Taubin
\cite{taubin} y se simplifica con métricas de error cuádrico \cite{quadric}.

\subsection{Estimación de pose con marcadores}
Los marcadores ArUco son patrones cuadrados con identificador único
\cite{aruco}. Conocida la posición de cada marcador en el plato, la pose de la
cámara se obtiene resolviendo el problema perspectiva-$n$-puntos con todos los
marcadores visibles, aun cuando el objeto tape algunos.

\subsection{Comunicación segura}
TLS~1.3 \cite{tls13} ofrece cifrado autenticado y secreto hacia adelante. Con
autenticación mutua (mTLS) cada extremo presenta un certificado firmado por una
autoridad propia. Se prefiere ChaCha20-Poly1305 \cite{chacha} porque la
Raspberry~Pi~3 no acelera AES por hardware.

% =====================================================================
\section{Ambiente de desarrollo}

\subsection{Hardware}
\begin{itemize}
  \item Computadora anfitriona (\emph{host}): captura, coordinación, fusión,
        mallado y visor. \pendiente{modelo, CPU, RAM y sistema operativo}
  \item Computadora \emph{worker B}: fusión parcial.
        \pendiente{modelo, CPU, RAM y sistema operativo}
  \item Raspberry~Pi~3 con Raspberry Pi OS Lite de 64 bits: controlador,
        holograma y fusión ligera. \pendiente{versión exacta del kernel}
  \item iPhone con cámara TrueDepth. \pendiente{modelo}
  \item Cinco pulsadores, plato giratorio manual de 180~mm, monitor HDMI y panel
        transparente a $45^{\circ}$.
\end{itemize}

\subsection{Software y bibliotecas}
La Tabla~\ref{tab:software} resume las dependencias principales.

\begin{table}[!t]
\caption{Herramientas y bibliotecas}
\label{tab:software}
\centering
\footnotesize
\begin{tabularx}{\columnwidth}{lX}
\toprule
\textbf{Herramienta} & \textbf{Uso}\\
\midrule
Python $\geq$ 3.11, uv & Servicios y manejo de dependencias\\
NumPy $\geq$ 1.26 & Núcleo TSDF propio\\
OpenCV $\geq$ 4.8 & Marcadores ArUco y estimación de pose\\
SciPy, scikit-image & Mallado y métricas\\
fast-simplification & Decimación cuádrica\\
msgpack, zstandard & Serialización y compresión\\
record3d $\geq$ 1.4 & Flujo RGB-D del iPhone por USB \cite{record3d}\\
C99, raylib 5.5, CMake & Visor (GLSL 330 y GLSL ES 100) \cite{raylib}\\
Kernel de Linux, dtc & Módulo \texttt{holo-io} y \emph{device tree}\\
OpenSSL & Autoridad certificadora y certificados\\
pytest $\geq$ 8 & Pruebas unitarias\\
Git y GitHub & Control de versiones y \emph{pull requests}\\
\bottomrule
\end{tabularx}
\end{table}

% =====================================================================
\section{Atributos: Aprendizaje continuo}

\subsection{Necesidades de aprendizaje}
\begin{itemize}
  \item \textbf{Conocimientos:} subsistema GPIO, interrupciones, temporizadores y
        \emph{device tree} en Linux; geometría proyectiva y pose de cámara;
        fusión volumétrica TSDF; TLS~1.3 e infraestructura de clave pública.
  \item \textbf{Habilidades:} depurar código de kernel, programar protocolos de
        red asíncronos, medir rendimiento y escribir \emph{shaders} para GLES2.
  \item \textbf{Destrezas:} cableado de pulsadores y construcción de la maqueta.
  \item \textbf{Actitudes:} atacar primero los riesgos técnicos, documentar las
        decisiones y coordinar el trabajo en equipo.
\end{itemize}

\subsection{Tecnologías posibles y seleccionadas}
La Tabla~\ref{tab:alternativas} compara las alternativas consideradas.

\begin{table}[!t]
\caption{Alternativas tecnológicas}
\label{tab:alternativas}
\centering
\footnotesize
\begin{tabularx}{\columnwidth}{XXX}
\toprule
\textbf{Alternativa} & \textbf{Elegida} & \textbf{Motivo}\\
\midrule
App iOS propia & Record3D & No se disponía de macOS para compilar\\
Registro ICP entre vistas & Anillo ArUco & ICP es frágil con el ruido del sensor\\
Transmitir video & Enviar estado de cámara & Menor latencia y carga en la Pi\\
Plano del plato con RANSAC & Plano dado por ArUco & Es conocido y determinista\\
Directorio compartido & Mensajes sobre mTLS & Lo exige el enunciado\\
Tornamesa motorizada & Giro manual (motor opcional) & Reduce riesgo y costo\\
\bottomrule
\end{tabularx}
\end{table}

\subsection{Acciones implementadas y evidencia}
\begin{enumerate}
  \item \textbf{Planificación por hitos:} nueve hitos (M0--M8) con tareas,
        dependencias y criterios de aceptación en
        \texttt{plan/milestones-and-actions.md}.
  \item \textbf{Pruebas de viabilidad:} captura por USB y pose con el anillo se
        validaron antes de construir sobre ellas (\texttt{tests/spikes/}).
  \item \textbf{Pruebas automatizadas:} pruebas unitarias en
        \texttt{tests/unit/} y una prueba sintética con el \emph{Stanford bunny}
        (cortesía del Stanford Computer Graphics Laboratory).
  \item \textbf{Control de versiones:} ramas por hito e integración mediante
        \emph{pull requests}; historial entre el 8 y el 24 de noviembre de 2026.
  \item \textbf{Búsqueda de información y repaso:}
        \pendiente{fuentes consultadas, reuniones, horas y responsables}
\end{enumerate}

\pendiente{insertar capturas como evidencia: historial de commits, pull requests,
salida de pytest, fotos del montaje, actas de reunión}

\subsection{Evaluación crítica}
La Tabla~\ref{tab:evaluacion} evalúa la eficiencia de las acciones.
\pendiente{el equipo debe revisar esta evaluación y ajustarla a su experiencia}

\begin{table}[!t]
\caption{Evaluación de las acciones}
\label{tab:evaluacion}
\centering
\footnotesize
\begin{tabularx}{\columnwidth}{lXX}
\toprule
\textbf{Acción} & \textbf{Resultado} & \textbf{Limitación}\\
\midrule
Hitos & Permiten trabajar frentes en paralelo & Sin fechas ni responsables asignados\\
Viabilidad & Confirmó captura y pose con hardware real & Quedaron pendientes las pruebas en la Pi\\
Pruebas & Verifican la fusión y el protocolo sin hardware & No sustituyen la prueba del controlador\\
Versionado & Trazabilidad de cada cambio & Actividad concentrada en pocos días\\
\bottomrule
\end{tabularx}
\end{table}

% =====================================================================
\section{Detalles del diseño}

\subsection{Arquitectura}
El sistema usa una topología en estrella (Fig.~\ref{fig:arquitectura}). El
\emph{host} escucha en el puerto TCP 7400; \emph{worker B} y la Pi abren la
conexión hacia él y no se comunican entre sí.

\begin{figure*}[!t]
\centering
\begin{tikzpicture}[node distance=5mm and 24mm]
  \node[caja] (cap) {capture-service};
  \node[caja, below=of cap] (coord) {coordinator\\astrea-host};
  \node[caja, below=of coord] (wh) {worker};
  \node[caja, below=of wh] (vis) {renderer\\(viewer)};

  \node[caja, left=of cap] (phone) {iPhone TrueDepth\\Record3D};

  \node[caja, left=34mm of coord] (wb) {worker};

  \node[caja, right=34mm of coord] (agent) {device-agent};
  \node[caja, below=of agent] (wp) {worker (ligero)};
  \node[caja, below=of wp] (holo) {renderer\\(hologram)};
  \node[caja, above=of agent] (drv) {holo-io\\/dev/holo-input};
  \node[caja, above=of drv] (btn) {5 pulsadores\\GPIO};
  \node[caja, right=8mm of holo] (pant) {Pepper's\\Ghost};

  \begin{scope}[on background layer]
    \node[grupo, fit=(cap)(coord)(wh)(vis), label=above:{\footnotesize\bfseries Nodo 1: Host}] (g1) {};
    \node[grupo, fit=(wb), label=above:{\footnotesize\bfseries Nodo 2: Worker B}] (g2) {};
    \node[grupo, fit=(drv)(agent)(wp)(holo), label={[xshift=18mm]above:{\footnotesize\bfseries Nodo 3: Pi 3}}] (g3) {};
  \end{scope}

  \draw[flecha] (phone) -- node[etiqueta]{USB} (cap);
  \draw[doble] (coord) -- node[etiqueta]{mTLS\\N1} (wb);
  \draw[doble] (coord) -- node[etiqueta]{mTLS\\N1, N2, N3} (agent);
  \draw[flecha] (btn) -- node[etiqueta, right]{IRQ} (drv);
  \draw[flecha] (drv) -- (agent);
  \draw[flecha] (agent.east) -- ++(5mm,0) |- node[etiqueta, pos=0.25]{socket\\Unix} (holo.east);
  \draw[flecha] (holo) -- (pant);
\end{tikzpicture}
\caption{Arquitectura. N1: trabajos de reconstrucción; N2: envío del modelo;
N3: entradas y estado de cámara. \texttt{astrea-host}, \texttt{device-agent} y el
socket del visor están diseñados pero aún no implementados.}
\label{fig:arquitectura}
\end{figure*}

\subsection{Componentes}
La Fig.~\ref{fig:componentes} muestra los paquetes de software y sus
dependencias.

\begin{figure*}[!t]
\centering
\begin{tikzpicture}[node distance=12mm and 8mm]
  \node[caja] (capt) {capture\_service\\Record3D, ArUco,\\keyframes};
  \node[caja, right=of capt] (coo) {coordinator\\reparto, merge,\\mallado, LOD};
  \node[caja, right=of coo] (wor) {worker\\fusión parcial};
  \node[caja, right=of wor] (dev) {device\_agent\\botones y modelo\\(pendiente)};
  \node[caja, below=18mm of coo, xshift=18mm, minimum width=105mm]
        (com) {astrea\_common: TSDF, preprocesado, protocolo, framing, TLS, configuración, anillo ArUco, GLB};
  \node[caja, below=of com, xshift=-30mm] (ren) {renderer\\C + raylib};
  \node[caja, below=of com, xshift=30mm] (kern) {holo-io\\módulo de kernel};

  \foreach \n in {capt,coo,wor,dev}
    \draw[flecha] (\n) -- (\n |- com.north);
  \draw[flecha, dashed] (ren) -- node[etiqueta]{carga GLB} (ren |- com.south);
  \draw[flecha, dashed] (kern) -- node[etiqueta]{HoloEvent} (kern |- com.south);
\end{tikzpicture}
\caption{Componentes. Las líneas continuas son dependencias de código; las
punteadas, contratos de formato compartidos entre lenguajes.}
\label{fig:componentes}
\end{figure*}

\subsection{Controlador \texttt{holo-io}}
Es un \emph{platform driver} asociado a un nodo de \emph{device tree} compatible
con \texttt{astreaview,holo-io}. Sus características son:
\begin{itemize}
  \item interrupción por ambos flancos en cada GPIO;
  \item \emph{debounce} con \texttt{hrtimer} (10~ms por defecto);
  \item repetición automática cada 50~ms en los botones de movimiento;
  \item cola \texttt{kfifo} de 256 eventos y \emph{wait queue}, con
        \texttt{read} y \texttt{poll};
  \item \texttt{ioctl} para ajustar tiempos y estadísticas en \texttt{sysfs};
  \item apertura exclusiva de \texttt{/dev/holo-input}.
\end{itemize}

Cada evento es un registro de 16 bytes:
\begin{lstlisting}[language=C]
struct HoloEvent {
    __u64 timestampNs;
    __u16 source;
    __u16 code;
    __s32 value;   /* 0 release, 1 press, 2 repeat */
};
\end{lstlisting}

La Tabla~\ref{tab:botones} muestra la asignación de pines. Cada pulsador va entre
el GPIO y tierra, con resistencia de \emph{pull-up} interna.

\begin{table}[!t]
\caption{Pulsadores (numeración BCM)}
\label{tab:botones}
\centering
\footnotesize
\begin{tabular}{ccl}
\toprule
\textbf{GPIO} & \textbf{Código} & \textbf{Acción}\\
\midrule
4 & 1 & Girar a la derecha\\
27 & 2 & Girar a la izquierda\\
22 & 3 & Acercar\\
23 & 4 & Alejar\\
24 & 5 & Activar o desactivar autorrotación\\
\bottomrule
\end{tabular}
\end{table}

\subsection{Captura y reconstrucción}
\textbf{Captura.} El anillo tiene 12 marcadores \texttt{DICT\_4X4\_50} de 24~mm,
con centros a 70~mm del eje. Se guarda una vista cuando el ángulo avanzó al menos
$10^{\circ}$ y la pose se mantiene estable cerca de 0.5~s, promediando los
cuadros de esa ventana.

\textbf{Preprocesado.} Por cada vista se recorta la profundidad entre 200 y
500~mm, se conserva un cilindro de 100~mm de radio sobre el plato y se eliminan
píxeles aislados y componentes pequeñas. Solo usa datos de la propia vista, de
modo que repartir las vistas no altera el resultado.

\textbf{Integración y salida.} Cada observación se pondera por ángulo de
incidencia y profundidad. El tamaño de vóxel por defecto es 0.65~mm. De la malla
se exportan tres niveles de detalle en GLB: completo, alto (unos 30\,000
triángulos) y bajo (unos 8\,000, para la Pi).

\subsection{Distribución}
El coordinador fija la rejilla, las poses y los parámetros, y reparte las vistas
por turnos (\emph{round robin}). Cada \emph{worker} fusiona en un hilo aparte,
para no detener los latidos de la conexión, y devuelve solo los bloques de $32^3$
vóxeles con peso positivo, comprimidos con zstd. El coordinador los combina con
la ecuación~(1). El reparto proporcional al rendimiento y la reasignación ante
fallos están pendientes: hoy la caída de un \emph{worker} aborta la
reconstrucción.

\subsection{Protocolo y seguridad}
Cada trama tiene la forma \texttt{[u32 length][u8 type][msgpack]}, con un máximo
de 4~MiB. La conexión usa exclusivamente TLS~1.3 con certificados ECDSA~P-256 de
una autoridad propia; el nombre del nodo del certificado debe coincidir con el
declarado en \texttt{HELLO}. Hay latidos cada 2~s y se cierra el enlace tras 6~s
de silencio. La Tabla~\ref{tab:mensajes} lista los mensajes.

\begin{table}[!t]
\caption{Mensajes del protocolo}
\label{tab:mensajes}
\centering
\footnotesize
\begin{tabularx}{\columnwidth}{lX}
\toprule
\textbf{Mensaje} & \textbf{Propósito}\\
\midrule
HELLO & Identificación y versión\\
HEARTBEAT & Detección de caídas y medición de ida y vuelta\\
JOB\_ASSIGN & Vistas, poses, rejilla y parámetros\\
JOB\_RESULT & Bloques del volumen parcial\\
INPUT\_EVENT & Evento de un botón\\
CAMERA\_STATE & Giro, inclinación fija, acercamiento y modelo\\
MODEL\_PUSH & Modelo en fragmentos con SHA-256\\
ROTATE, ROTATE\_DONE & Reservados para la tornamesa motorizada\\
\bottomrule
\end{tabularx}
\end{table}

\subsection{Diagramas de secuencia}
Las Figs.~\ref{fig:sec-escaneo} y~\ref{fig:sec-botones} muestran los dos flujos
principales.

\begin{figure*}[!t]
\centering
\begin{tikzpicture}[x=1cm, y=0.8cm]
  \foreach \x/\n in {0/Usuario, 3/Captura, 6.3/Coordinator, 9.9/Workers (3), 13.5/Visores} {
    \node[caja] at (\x,0) {\n};
    \draw[dashed, gray] (\x,-0.5) -- (\x,-9.4);
  }
  \draw[flecha] (0,-1)   -- node[etiqueta, above]{gira el plato} (3,-1);
  \draw[flecha] (3,-1.7) -- ++(0.8,0) -- ++(0,-0.6) node[etiqueta, midway, right]{pose y keyframe} -- (3,-2.3);
  \draw[flecha] (3,-3.2) -- node[etiqueta, above]{sesión de vistas} (6.3,-3.2);
  \draw[flecha] (6.3,-4.2) -- node[etiqueta, above]{JOB\_ASSIGN} (9.9,-4.2);
  \draw[flecha] (9.9,-5) -- ++(0.8,0) -- ++(0,-0.6) node[etiqueta, midway, right]{fusión TSDF} -- (9.9,-5.6);
  \draw[flecha] (9.9,-6.4) -- node[etiqueta, above]{JOB\_RESULT} (6.3,-6.4);
  \draw[flecha] (6.3,-7.1) -- ++(0.8,0) -- ++(0,-0.6) node[etiqueta, midway, right]{merge y malla} -- (6.3,-7.7);
  \draw[flecha] (6.3,-8.7) -- node[etiqueta, above]{GLB / MODEL\_PUSH} (13.5,-8.7);
\end{tikzpicture}
\caption{Secuencia de escaneo. El envío del modelo a la Pi corresponde al diseño
y está pendiente.}
\label{fig:sec-escaneo}
\end{figure*}

\begin{figure*}[!t]
\centering
\begin{tikzpicture}[x=1cm, y=0.8cm]
  \foreach \x/\n in {0/Pulsador, 2.8/holo-io, 6/device-agent, 9.6/astrea-host, 13.2/Visores} {
    \node[caja] at (\x,0) {\n};
    \draw[dashed, gray] (\x,-0.5) -- (\x,-6.2);
  }
  \draw[flecha] (0,-1)   -- node[etiqueta, above]{flanco (IRQ)} (2.8,-1);
  \draw[flecha] (2.8,-1.7) -- ++(0.8,0) -- ++(0,-0.6) node[etiqueta, midway, right]{debounce} -- (2.8,-2.3);
  \draw[flecha] (2.8,-3) -- node[etiqueta, above]{HoloEvent} (6,-3);
  \draw[flecha] (6,-3.7) -- node[etiqueta, above]{INPUT\_EVENT} (9.6,-3.7);
  \draw[flecha] (9.6,-4.6) -- node[etiqueta, above]{CAMERA\_STATE (a la Pi)} (6,-4.6);
  \draw[flecha] (9.6,-5.5) -- node[etiqueta, above]{estado local} (13.2,-5.5);
\end{tikzpicture}
\caption{Secuencia de botones a holograma (diseño). El visor de la Pi recibe el
estado a través del \texttt{device-agent}.}
\label{fig:sec-botones}
\end{figure*}

\subsection{Funcionalidades}
La Fig.~\ref{fig:funciones} agrupa las funciones del sistema.

\begin{figure}[!t]
\centering
\begin{tikzpicture}[node distance=3mm,
  rama/.style={caja, text width=24mm},
  hoja/.style={caja, text width=46mm, align=left}]
  \node[rama] (a) {Captura};
  \node[hoja, right=5mm of a] (a1) {Flujo RGB-D, pose ArUco, keyframes automáticos, indicador de cobertura};
  \node[rama, below=9mm of a] (b) {Reconstrucción};
  \node[hoja, right=5mm of b] (b1) {Preprocesado, TSDF, malla, niveles de detalle};
  \node[rama, below=9mm of b] (c) {Distribución};
  \node[hoja, right=5mm of c] (c1) {Reparto, fusión parcial, merge exacto, reasignación*};
  \node[rama, below=9mm of c] (d) {Seguridad};
  \node[hoja, right=5mm of d] (d1) {mTLS, certificados por nodo, validación de mensajes};
  \node[rama, below=9mm of d] (e) {Interacción*};
  \node[hoja, right=5mm of e] (e1) {Botones, teclado, autorrotación, visor y holograma sincronizados};
  \foreach \n in {a,b,c,d,e} \draw[flecha] (\n) -- (\n1);
\end{tikzpicture}
\caption{Funcionalidades. El asterisco indica funciones aún incompletas.}
\label{fig:funciones}
\end{figure}

% =====================================================================
\section{Resultados}

\subsection{Captura}
El flujo por USB entregó cerca de 27 cuadros por segundo a $480\times640$ con sus
parámetros intrínsecos. La pose con el anillo real se mantuvo estable en reposo,
en pasos de $30^{\circ}$ y con oclusión parcial.
\pendiente{error angular medido contra marcas en el plato; meta: menor que 1 grado}

\subsection{Reconstrucción}
En la prueba sintética sin ruido, la distancia de la reconstrucción al modelo
original tuvo media de 0.16 vóxeles y percentil 95 de 0.38 vóxeles; el umbral
era 2 y 4 vóxeles. El nivel alto pesa 594~KiB (485~KiB con zstd) y el bajo
160~KiB (131~KiB). \pendiente{resultados con capturas reales e imágenes}

\subsection{Distribución}
Con un escaneo real de 59 vistas y tres \emph{workers} en una misma máquina, la
rejilla de $332\times334\times256$ (568~MiB) tomó 110~s en total: cerca de 100~s
de fusión por \emph{worker} y 3.9~s de \emph{merge}. Cada proceso llegó a unos
1.3~GB de memoria. Con vóxel de 1.3~mm, la corrida distribuida y la de un nodo
produjeron 54\,967 caras en el nivel completo.
\pendiente{tabla de speedup 1 vs 3 workers en las tres máquinas reales}

\subsection{Red segura}
Los tres nodos reales se conectaron e intercambiaron latidos.
\pendiente{captura de Wireshark y prueba de rechazo de certificados inválidos}

\subsection{Controlador, visor y latencia}
El módulo compila sin advertencias contra los encabezados 6.12 de x86, pero aún
no se ha cargado en la Raspberry~Pi.
\pendiente{salida de dmesg, estadísticas de sysfs, prueba de 100 pulsaciones,
cuadros por segundo en la Pi, latencia del botón al holograma, fotos}

% =====================================================================
\section{Instrucciones de utilización}

\subsection{Requisitos e instalación}
\begin{lstlisting}[language=bash]
git clone <URL del repositorio> astrea-view
cd astrea-view
uv sync --extra record3d --extra reconstruction
uv run pytest
\end{lstlisting}

En Linux, la captura por USB requiere \texttt{usbmuxd} y la función de
transmisión USB de Record3D.

\subsection{Certificados y configuración}
\begin{lstlisting}[language=bash]
pki/make-certs.sh --out pki/out --ip host=<IP>
\end{lstlisting}
Copie a cada nodo \texttt{ca.crt}, su certificado y su clave. En
\texttt{config.toml} ajuste \texttt{hub.address}, \texttt{hub.bind-address} y las
rutas de \texttt{[tls]} para usar \texttt{pki/out/}.

\subsection{Anillo y captura}
\begin{lstlisting}[language=bash]
uv run hardware/aruco-ring/generate-ring.py \
    --out-dir hardware/aruco-ring/out
uv run astrea-capture --name taza
\end{lstlisting}
Imprima el anillo al 100\,\% de escala, coloque el objeto al centro y gire el
plato con pausas.

\subsection{Reconstrucción}
En un solo nodo:
\begin{lstlisting}[language=bash]
uv run astrea-reconstruct out/scans/taza
\end{lstlisting}
Distribuida (primero el host y luego cada worker):
\begin{lstlisting}[language=bash]
# host
uv run astrea-coordinator --config config.toml \
    reconstruct out/scans/taza --min-workers 3 --voxel-mm 1.3
# host y worker B
uv run astrea-worker --name host
uv run astrea-worker --name worker-b
# Raspberry Pi
OPENSSL_CONF=pki/openssl-chacha20.cnf \
    uv run astrea-worker --name pi
\end{lstlisting}

\subsection{Visor}
\begin{lstlisting}[language=bash]
cmake -S renderer -B renderer/build -DCMAKE_BUILD_TYPE=Release
# en la Pi agregar: -DRENDERER_PLATFORM=DRM
cmake --build renderer/build -j
renderer/build/renderer --mode viewer \
    --model out/scans/taza/model/lod-high.glb --up-axis z
\end{lstlisting}

\subsection{Controlador (Raspberry Pi)}
\begin{lstlisting}[language=bash]
cd driver/holo-io
make
sudo make install
echo "dtoverlay=holo-io" | sudo tee -a /boot/firmware/config.txt
sudo reboot
# despues de reiniciar, en la misma carpeta
./holo-reader
./check-presses.sh 100
\end{lstlisting}

% =====================================================================
\section{Conclusiones}
\pendiente{revisar al terminar las pruebas en hardware}
\begin{itemize}
  \item La fusión TSDF por promedio ponderado permite distribuir las vistas y
        obtener un resultado equivalente al de un solo nodo.
  \item Un anillo de marcadores sustituye a un codificador o motor para conocer
        la pose en un plato girado a mano.
  \item La memoria, más que el cálculo, limita la participación de la Raspberry
        Pi: obliga a usar una rejilla más gruesa.
  \item Separar el preprocesado por vista y fijar la rejilla en el coordinador es
        lo que hace reproducible la distribución.
\end{itemize}

\section{Sugerencias y recomendaciones}
\begin{itemize}
  \item Probar en el hardware definitivo desde las primeras semanas, en especial
        el módulo de kernel.
  \item Reasignar el trabajo de un nodo caído y repartir según el rendimiento de
        cada uno.
  \item Agregar la tornamesa motorizada para reforzar el componente de hardware.
  \item Usar certificados propios en la demostración, nunca los de ejemplo.
\end{itemize}

% =====================================================================
\begin{thebibliography}{00}
\bibitem{ldd3} J. Corbet, A. Rubini y G. Kroah-Hartman, \emph{Linux Device
Drivers}, 3.\textsuperscript{a} ed. O'Reilly Media, 2005.
\bibitem{curless} B. Curless y M. Levoy, ``A volumetric method for building
complex models from range images,'' en \emph{Proc. SIGGRAPH}, 1996, pp. 303--312.
\bibitem{kinectfusion} R. A. Newcombe \emph{et al.}, ``KinectFusion: Real-time
dense surface mapping and tracking,'' en \emph{Proc. IEEE ISMAR}, 2011,
pp. 127--136.
\bibitem{marching} W. E. Lorensen y H. E. Cline, ``Marching cubes: A high
resolution 3D surface construction algorithm,'' en \emph{Proc. SIGGRAPH}, 1987,
pp. 163--169.
\bibitem{taubin} G. Taubin, ``A signal processing approach to fair surface
design,'' en \emph{Proc. SIGGRAPH}, 1995, pp. 351--358.
\bibitem{quadric} M. Garland y P. S. Heckbert, ``Surface simplification using
quadric error metrics,'' en \emph{Proc. SIGGRAPH}, 1997, pp. 209--216.
\bibitem{aruco} S. Garrido-Jurado, R. Mu\~noz-Salinas, F. J. Madrid-Cuevas y
M. J. Mar\'in-Jim\'enez, ``Automatic generation and detection of highly reliable
fiducial markers under occlusion,'' \emph{Pattern Recognition}, vol. 47, n.o 6,
pp. 2280--2292, 2014.
\bibitem{tls13} E. Rescorla, ``The Transport Layer Security (TLS) Protocol
Version 1.3,'' RFC 8446, IETF, 2018.
\bibitem{chacha} Y. Nir y A. Langley, ``ChaCha20 and Poly1305 for IETF
Protocols,'' RFC 8439, IETF, 2018.
\bibitem{raylib} R. Santamaria, ``raylib.'' [En línea]. Disponible:
\url{https://www.raylib.com}
\bibitem{record3d} M. Simonik, ``Record3D.'' [En línea]. Disponible:
\url{https://record3d.app}
\end{thebibliography}

\end{document}
